%---------------------------------------------------------------------
% The fourteen questions, each with its past-paper provenance.
%---------------------------------------------------------------------

\section*{خبر: مفہوم، اقدار اور اقسام}
\markboth{خبر: مفہوم، اقدار، اقسام}{}

\begin{sawal}{سوال \textenglish{1} \ \textnormal{\small(بہار \textenglish{2013}، سوال \textenglish{1} اور خزاں \textenglish{2013}، سوال \textenglish{1})}}
خبر کا لغوی اور اصطلاحی مفہوم بیان کریں، اور ابلاغِ عامہ میں خبر کے کردار پر
روشنی ڈالیں۔
\end{sawal}

\begin{hal}
\textbf{لغوی مفہوم۔} ''خبر'' عربی لفظ ہے جس کے معنی ہیں اطلاع، آگاہی یا کسی
بات کی اطلاع دینا۔ انگریزی لفظ \textenglish{News} کے بارے میں مشہور ہے کہ یہ
چاروں سمتوں --- \textenglish{North, East, West, South} --- کے پہلے حروف سے بنا
ہے، یعنی ہر طرف سے آنے والی اطلاعات۔ (یہ توجیہ عام ہے، اگرچہ لغوی طور پر
\textenglish{News} دراصل \textenglish{new} کی جمع ہے۔)

\medskip
\textbf{اصطلاحی مفہوم۔} خبر کسی ایسے تازہ، سچے اور غیر معمولی واقعے یا خیال کی
بروقت اطلاع ہے جس میں قارئین کی ایک بڑی تعداد دلچسپی رکھتی ہو اور جو ان پر اثر
انداز ہوتا ہو۔

\smallskip
مشہور تعریفیں:
\begin{itemize}\setlength{\itemsep}{1pt}
  \item ''کتا آدمی کو کاٹے تو خبر نہیں، آدمی کتے کو کاٹے تو خبر ہے۔'' ---
        یعنی خبر کی جان \emph{غیر معمولی پن} ہے۔
  \item ''خبر وہ ہے جسے کوئی چھپانا چاہتا ہو؛ باقی سب اشتہار ہے۔''
\end{itemize}

\medskip
\textbf{ابلاغِ عامہ میں خبر کا کردار:}
\begin{itemize}\setlength{\itemsep}{1pt}
  \item \textbf{اطلاع} --- عوام کو حالات سے باخبر رکھنا؛ یہی ابلاغ کا بنیادی
        کام ہے۔
  \item \textbf{نگرانی} --- خبر ہی اداروں کے احتساب کا ذریعہ بنتی ہے۔
  \item \textbf{رائے عامہ کی تشکیل} --- لوگ خبر کی بنیاد پر رائے بناتے ہیں۔
  \item \textbf{فیصلہ سازی} --- شہری، تاجر اور حکومت سب خبر پر انحصار کرتے ہیں۔
  \item \textbf{ربط} --- معاشرے کے مختلف طبقات کو ایک دوسرے سے جوڑنا۔
  \item \textbf{معاشی بنیاد} --- خبر ہی قاری کو کھینچتی ہے، اور قاری اشتہار
        لاتا ہے، جس پر پورا ادارہ کھڑا ہے۔
\end{itemize}

\jawab{لغوی طور پر خبر کے معنی اطلاع و آگاہی ہیں۔ اصطلاحاً خبر کسی تازہ، سچے
اور غیر معمولی واقعے کی بروقت اطلاع ہے جو بڑی تعداد میں لوگوں کے لیے دلچسپ یا
اثر انگیز ہو۔ ابلاغِ عامہ میں خبر اطلاع، نگرانی، رائے عامہ کی تشکیل، فیصلہ سازی
اور معاشرتی ربط کا کام کرتی ہے۔}
\end{hal}

\begin{sawal}{سوال \textenglish{2} \ \textnormal{\small(بہار \textenglish{2013}، سوال \textenglish{1})}}
خبر کی اقدار پر روشنی ڈالیں۔
\end{sawal}

\begin{hal}
\textbf{خبر کی اقدار} (\textenglish{News Values}) وہ کسوٹیاں ہیں جن پر پرکھ کر
ایڈیٹر فیصلہ کرتا ہے کہ کون سی اطلاع خبر بننے کے قابل ہے اور کتنی نمایاں جگہ
پائے گی۔ جتنی زیادہ اقدار موجود ہوں، خبر اتنی ہی بڑی۔

\medskip
\begin{center}
\begin{tabular}{@{}p{3.2cm} p{5.4cm} p{5.8cm}@{}}
\toprule
\textbf{قدر} & \textbf{مطلب} & \textbf{مثال} \\
\midrule
تازگی & واقعہ جتنا نیا، خبر اتنی بڑی & آج کا اعلان بمقابلہ ہفتے پرانا \\
قربت & جگہ یا دلچسپی کے اعتبار سے قریب & اپنے شہر کا حادثہ زیادہ اہم \\
اہمیت و اثر & کتنے لوگ متاثر ہوں گے & بجلی کی قیمت میں اضافہ \\
نمایاں شخصیت & مشہور فرد سے متعلق & وزیرِ اعظم یا کھلاڑی کا بیان \\
تصادم & اختلاف، مقابلہ، جنگ & سیاسی محاذ آرائی، کھیل کا فائنل \\
غیر معمولی پن & جو روز نہ ہوتا ہو & انوکھا واقعہ یا ریکارڈ \\
انسانی دلچسپی & جذبات کو چھونے والی بات & کسی غریب بچے کی کامیابی \\
نتائج & آگے کیا ہوگا & بجٹ کے معیشت پر اثرات \\
تعداد و حجم & جتنا بڑا، اتنا اہم & دس ہلاکتیں بمقابلہ ایک \\
\bottomrule
\end{tabular}
\end{center}

\jawab{خبر کی اقدار وہ معیار ہیں جن سے خبر کی اہمیت جانچی جاتی ہے: تازگی،
قربت، اہمیت و اثر، نمایاں شخصیت، تصادم، غیر معمولی پن، انسانی دلچسپی، نتائج
اور تعداد۔ جتنی زیادہ اقدار جمع ہوں، خبر اتنی نمایاں چھپتی ہے۔}
\end{hal}

\begin{sawal}{سوال \textenglish{3} \ \textnormal{\small(بہار \textenglish{2013}، سوال \textenglish{2})}}
موضوع کے لحاظ سے خبر کی اقسام بیان کریں۔
\end{sawal}

\begin{hal}
\begin{center}
\begin{tabular}{@{}p{3.6cm} p{10.8cm}@{}}
\toprule
سیاسی خبر & حکومت، جماعتیں، اسمبلی، انتخابات --- اخبار کی صفحۂ اول کی جان \\
جرائم کی خبر & قتل، ڈکیتی، منشیات؛ محتاط زبان اور قانونی احتیاط ضروری \\
عدالتی خبر & مقدمات اور فیصلے؛ زیرِ سماعت مقدمے پر تبصرہ توہینِ عدالت بن سکتا ہے \\
معاشی و کاروباری & بجٹ، مہنگائی، اسٹاک مارکیٹ، زرِ مبادلہ \\
کھیلوں کی خبر & میچ، کھلاڑی، ٹورنامنٹ؛ الگ صفحہ اور الگ اسلوب \\
تعلیمی خبر & نتائج، داخلے، پالیسی، امتحانات \\
صحت کی خبر & وبائیں، ہسپتال، طبی تحقیق \\
سائنس و ٹیکنالوجی & ایجادات، خلائی تحقیق، انٹرنیٹ \\
ثقافتی و شوبز & ادب، فنون، فلم، ٹی وی \\
حادثات و آفات & ٹریفک حادثہ، زلزلہ، سیلاب \\
بین الاقوامی خبر & دیگر ممالک کے حالات؛ عموماً خبر رساں اداروں سے \\
انسانی دلچسپی & جذباتی یا انوکھی کہانیاں \\
موسم کی خبر & روزمرہ کی سب سے زیادہ پڑھی جانے والی خبروں میں سے \\
\bottomrule
\end{tabular}
\end{center}

\jawab{موضوع کے اعتبار سے خبر کی اہم اقسام: سیاسی، جرائم، عدالتی، معاشی،
کھیل، تعلیمی، صحت، سائنس و ٹیکنالوجی، ثقافتی و شوبز، حادثات، بین الاقوامی،
انسانی دلچسپی اور موسم۔}
\end{hal}

\begin{sawal}{سوال \textenglish{4} \ \textnormal{\small(بہار \textenglish{2013}، سوال \textenglish{3})}}
خبر نگاری کے اصول اور اسلوبِ تحریر مثالوں سے واضح کریں۔
\end{sawal}

\begin{hal}
\textbf{خبر کے اجزا:}
\begin{center}
\begin{tabular}{@{}p{3.0cm} p{11.4cm}@{}}
\toprule
سرخی & خبر کا خلاصہ چند الفاظ میں؛ قاری کو کھینچنے کا کام \\
ڈیٹ لائن & جگہ اور تاریخ --- ''اسلام آباد، \textenglish{5} ستمبر'' \\
انٹرو / سرخبر & پہلا پیرا؛ سب سے اہم بات، عموماً \textenglish{25--35} الفاظ \\
باڈی & تفصیل، پس منظر، بیانات --- اہمیت کے گھٹتے ہوئے ترتیب سے \\
\bottomrule
\end{tabular}
\end{center}

\medskip
\textbf{بنیادی اصول --- چھ سوالات۔} ہر مکمل خبر ان کا جواب دیتی ہے:
\textbf{کیا} ہوا، \textbf{کون} ملوث تھا، \textbf{کب} ہوا، \textbf{کہاں} ہوا،
\textbf{کیوں} ہوا، اور \textbf{کیسے} ہوا۔ (انگریزی میں
\textenglish{5W + 1H}۔)

\medskip
\textbf{الٹا اہرام} (\textenglish{Inverted Pyramid})۔ خبر کا معیاری ڈھانچہ:
سب سے اہم بات سب سے اوپر، پھر اہمیت کے گھٹتے درجے میں تفصیل۔ فائدہ یہ کہ قاری
پہلی سطر ہی سے اصل بات جان لیتا ہے، اور سب ایڈیٹر جگہ کم ہو تو خبر نیچے سے کاٹ
سکتا ہے۔

\medskip
\textbf{اسلوبِ تحریر کے اصول:}
\begin{itemize}\setlength{\itemsep}{1pt}
  \item زبان \textbf{سادہ، رواں اور عام فہم} ہو؛ ادبی بھاری پن سے گریز۔
  \item جملے \textbf{مختصر} ہوں --- ایک جملے میں ایک بات۔
  \item \textbf{معلوم فعل} (\textenglish{active voice}) استعمال کیجیے:
        ''پولیس نے ملزم گرفتار کر لیا'' نہ کہ ''ملزم گرفتار کر لیا گیا''۔
  \item \textbf{غیر جانبداری} --- خبر میں اپنی رائے شامل نہ کیجیے۔
  \item ہر دعوے کا \textbf{حوالہ} دیجیے: ''ترجمان کے مطابق''۔
  \item اعداد، نام اور عہدے \textbf{بالکل درست} ہوں۔
  \item مبالغہ، القاب اور جذباتی الفاظ سے پرہیز۔
\end{itemize}

\jawab{خبر کے اجزا سرخی، ڈیٹ لائن، انٹرو اور باڈی ہیں۔ ہر خبر چھ سوالات (کیا،
کون، کب، کہاں، کیوں، کیسے) کا جواب دیتی ہے اور الٹے اہرام کے ڈھانچے میں لکھی
جاتی ہے۔ اسلوب سادہ، مختصر جملوں پر مبنی، معلوم فعل والا، غیر جانبدار اور
باحوالہ ہونا چاہیے۔}
\end{hal}

\begin{ghalti}
خبر اور اداریہ ایک نہیں۔ خبر میں \emph{اپنی رائے} لکھنا سب سے بڑی غلطی ہے۔
اگر کوئی رائے دینی ہو تو وہ کسی ذمہ دار شخص کے حوالے سے ہونی چاہیے ---
''وزیرِ اعظم نے کہا کہ\ldots''۔
\end{ghalti}

%---------------------------------------------------------------------
\section*{رپورٹر اور اس کا کام}
\markboth{رپورٹر اور اس کا کام}{}

\begin{sawal}{سوال \textenglish{5} \ \textnormal{\small(خزاں \textenglish{2013}، سوال \textenglish{3})}}
خبروں کے حصول کے ذرائع کون سے ہیں؟ تفصیل سے لکھیں۔
\end{sawal}

\begin{hal}
\textbf{(۱) ادارے کے اپنے ذرائع:} رپورٹر، نامہ نگار، علاقائی نمائندے اور
فوٹو گرافر --- سب سے قابلِ اعتماد ذریعہ، کیونکہ خبر براہِ راست ادارے کے آدمی
کی ہوتی ہے۔

\textbf{(۲) خبر رساں ادارے:} \textenglish{APP}، \textenglish{PPI}،
\textenglish{Reuters}، \textenglish{AFP} وغیرہ --- خاص طور پر بین الاقوامی اور
دور دراز کی خبروں کا بنیادی ذریعہ۔

\textbf{(۳) سرکاری ذرائع:} پریس ریلیز، پریس انفارمیشن ڈیپارٹمنٹ
(\textenglish{PID})، وزارتوں کے ترجمان، سرکاری اعلامیے۔

\textbf{(۴) پریس کانفرنس اور بریفنگ:} جہاں کئی صحافی بیک وقت سوال کر سکتے ہیں۔

\textbf{(۵) ذاتی مشاہدہ:} رپورٹر خود موقع پر موجود ہو --- سب سے مستند صورت۔

\textbf{(۶) انٹرویو:} متعلقہ شخصیت سے براہِ راست گفتگو۔

\textbf{(۷) دستاویزات و ریکارڈ:} عدالتی فیصلے، بجٹ، رپورٹیں، اعداد و شمار ---
تحقیقی خبر کی ریڑھ کی ہڈی۔

\textbf{(۸) مخبر اور ذاتی روابط:} تھانہ، ہسپتال، دفاتر میں تعلقات؛ ان کے نام
صیغۂ راز میں رکھنا صحافی کا اخلاقی فرض ہے۔

\textbf{(۹) عوامی اجتماعات:} اسمبلی کی کارروائی، جلسے، احتجاج، عدالتیں۔

\textbf{(۱۰) جدید ذرائع:} انٹرنیٹ، سرکاری ویب سائٹس، سوشل میڈیا --- تیز مگر
سب سے کم قابلِ اعتماد؛ تصدیق کے بغیر استعمال خطرناک ہے۔

\jawab{اہم ذرائع: ادارے کے اپنے رپورٹر و نامہ نگار، خبر رساں ادارے، سرکاری
ذرائع و پریس ریلیز، پریس کانفرنس، ذاتی مشاہدہ، انٹرویو، دستاویزات، مخبر و ذاتی
روابط، عوامی اجتماعات اور انٹرنیٹ و سوشل میڈیا۔ جدید ذرائع تیز ہیں مگر تصدیق
لازمی ہے۔}
\end{hal}

\begin{sawal}{سوال \textenglish{6} \ \textnormal{\small(بہار \textenglish{2013}، سوال \textenglish{4} اور خزاں \textenglish{2013}، سوال \textenglish{4})}}
رپورٹر کسے کہتے ہیں؟ اس کے فرائض اور اوصاف پر روشنی ڈالیں۔
\end{sawal}

\begin{hal}
یہ موضوع دونوں پرچوں میں آیا --- سب سے اہم سوال یہی ہے۔

\medskip
\textbf{تعریف۔} رپورٹر وہ صحافی ہے جو خبر کے موقع تک پہنچتا ہے، حقائق جمع کرتا
ہے، ان کی تصدیق کرتا ہے اور خبر لکھ کر ادارے کو بھیجتا ہے۔ وہ اخبار کی
\emph{آنکھ اور کان} ہے۔

\medskip
\textbf{فرائض:}
\begin{itemize}\setlength{\itemsep}{1pt}
  \item اپنی بیٹ کی خبروں کا تعاقب اور موقع پر پہنچنا۔
  \item حقائق، اعداد و شمار اور بیانات جمع کرنا اور \textbf{تصدیق} کرنا۔
  \item متعلقہ افراد سے انٹرویو اور دونوں فریقوں کا مؤقف لینا۔
  \item خبر کو صحافتی اسلوب میں لکھ کر \textbf{ڈیڈ لائن} سے پہلے بھیجنا۔
  \item خبر کے پس منظر اور تسلسل پر نظر رکھنا (\textenglish{follow-up})۔
  \item ذرائع کے ساتھ تعلقات قائم رکھنا اور ان کا راز محفوظ رکھنا۔
  \item ضابطۂ اخلاق اور صحافتی قوانین کی پابندی۔
\end{itemize}

\textbf{اوصاف:}
\begin{itemize}\setlength{\itemsep}{1pt}
  \item \textbf{تجسس اور مشاہدہ} --- سوال اٹھانے کی عادت۔
  \item \textbf{دیانت اور غیر جانبداری} --- ذاتی پسند ناپسند سے بالاتر۔
  \item \textbf{زبان پر عبور} اور تیز رفتار تحریر۔
  \item \textbf{وسیع معلومات} --- سیاست، قانون، معیشت کی سمجھ۔
  \item \textbf{جفاکشی} --- ہر موسم اور ہر وقت کام۔
  \item \textbf{خوش اخلاقی و روابط} --- لوگ اسی سے بات کرتے ہیں جس پر بھروسہ ہو۔
  \item \textbf{جرأت} اور دباؤ میں غیر جانبدار رہنے کی ہمت۔
  \item \textbf{وقت کی پابندی} --- خبر وقت پر نہ پہنچے تو بے کار ہے۔
\end{itemize}

\jawab{رپورٹر وہ صحافی ہے جو موقع سے حقائق جمع کر کے، تصدیق کے بعد خبر لکھتا
ہے۔ فرائض: خبر کا تعاقب، تصدیق، انٹرویو، دونوں مؤقف، بروقت ترسیل اور
\textenglish{follow-up}۔ اوصاف: تجسس، دیانت، غیر جانبداری، زبان پر عبور، وسیع
معلومات، جفاکشی، روابط، جرأت اور وقت کی پابندی۔}
\end{hal}

\begin{sawal}{سوال \textenglish{7} \ \textnormal{\small(بہار \textenglish{2013}، سوال \textenglish{4} اور خزاں \textenglish{2013}، سوال \textenglish{7})}}
رپورٹر، نامہ نگار اور نمائندہ کے باہمی فرق کی وضاحت کریں۔
\end{sawal}

\begin{hal}
تینوں خبر بھیجتے ہیں، مگر \emph{مقام، حیثیت اور دائرۂ کار} مختلف ہے۔

\medskip
\begin{center}
\begin{tabular}{@{}p{2.4cm} p{4.0cm} p{3.8cm} p{3.6cm}@{}}
\toprule
 & \textbf{رپورٹر} & \textbf{نامہ نگار} & \textbf{نمائندہ} \\
\midrule
مقام & اخبار کے مرکزی دفتر والے شہر میں & کسی دوسرے شہر یا ملک میں & چھوٹے شہر، قصبے یا ضلع میں \\
تقرری & کل وقتی، تنخواہ دار & کل وقتی یا معاہدے پر & اکثر جز وقتی یا اعزازی \\
دائرۂ کار & مقررہ بیٹ (کرائم، عدالت، سیاست) & پورا شہر یا ملک & اپنا علاقہ \\
خبر کی نوعیت & روزانہ کی خبریں & مکتوب، تجزیہ، اہم خبریں & مقامی خبریں \\
اضافی کام & صرف صحافتی & صرف صحافتی & بعض اوقات اشاعت و اشتہار بھی \\
\bottomrule
\end{tabular}
\end{center}

\medskip
\textbf{یاد رکھنے کا آسان اصول:} رپورٹر \emph{مرکز} میں، نامہ نگار
\emph{دوسرے شہر} میں، اور نمائندہ \emph{مقامی سطح} پر۔

\jawab{رپورٹر مرکزی دفتر کے شہر میں مقررہ بیٹ پر کام کرنے والا کل وقتی صحافی
ہے؛ نامہ نگار دوسرے شہر یا ملک سے خبریں اور مکتوب بھیجتا ہے؛ اور نمائندہ چھوٹے
علاقے کا اکثر جز وقتی نمائندہ ہوتا ہے جو مقامی خبریں بھیجتا ہے۔}
\end{hal}

\begin{sawal}{سوال \textenglish{8} \ \textnormal{\small(بہار \textenglish{2013}، سوال \textenglish{6})}}
مختلف بیٹس سے تعلق رکھنے والے رپورٹروں کے فرائض اور صلاحیتوں پر تبصرہ کریں۔
\end{sawal}

\begin{hal}
\textbf{بیٹ کیا ہے؟} ''بیٹ'' (\textenglish{Beat}) اس مخصوص شعبے یا حلقے کو
کہتے ہیں جو کسی رپورٹر کے ذمے ہو۔ بیٹ کا نظام اس لیے بنا کہ رپورٹر ایک شعبے
میں \emph{مہارت} اور \emph{روابط} پیدا کر سکے۔

\medskip
\begin{center}
\begin{tabular}{@{}p{3.0cm} p{5.4cm} p{6.0cm}@{}}
\toprule
\textbf{بیٹ} & \textbf{فرائض} & \textbf{درکار صلاحیت} \\
\midrule
کرائم & تھانہ، جرائم، پولیس کارروائی & پولیس نظام کی سمجھ، محتاط زبان، ذاتی روابط \\
عدالتی & مقدمات اور فیصلے & قانونی اصطلاحات، توہینِ عدالت سے بچاؤ \\
سیاسی / اسمبلی & جماعتیں، اسمبلی، انتخابات & آئین و قواعد کا علم، مکمل غیر جانبداری \\
معاشی & بجٹ، مہنگائی، مارکیٹ & معاشی اصطلاحات، اعداد و شمار کی سمجھ \\
تعلیمی & امتحانات، پالیسی، ادارے & نظامِ تعلیم کی واقفیت \\
صحت & ہسپتال، وبائیں، طبی تحقیق & طبی اصطلاحات، خوف نہ پھیلانے کا شعور \\
کھیل & میچ، کھلاڑی، ٹورنامنٹ & کھیل کے قواعد، تیز رفتار رپورٹنگ \\
بلدیاتی & صفائی، پانی، سڑکیں & مقامی مسائل سے گہری واقفیت \\
سفارتی & سفارت خانے، دورے & بین الاقوامی تعلقات کا شعور، احتیاط \\
\bottomrule
\end{tabular}
\end{center}

\medskip
\textbf{مشترک صلاحیتیں:} ہر بیٹ کے رپورٹر میں تصدیق کی عادت، غیر جانبداری،
اصطلاحات پر عبور اور ذرائع کا مضبوط جال ضروری ہے۔

\jawab{بیٹ رپورٹر کے ذمے مخصوص شعبہ ہوتا ہے تاکہ وہ مہارت اور روابط پیدا کرے۔
کرائم، عدالتی، سیاسی، معاشی، تعلیمی، صحت، کھیل، بلدیاتی اور سفارتی بیٹ اہم
ہیں، اور ہر ایک کے لیے متعلقہ شعبے کی اصطلاحات، قواعد اور احتیاطوں کا علم
ضروری ہے۔}
\end{hal}

\begin{sawal}{سوال \textenglish{9} \ \textnormal{\small(خزاں \textenglish{2013}، سوال \textenglish{7})}}
علاقائی نامہ نگاری کے اصول بیان کریں۔
\end{sawal}

\begin{hal}
\textbf{تعریف۔} علاقائی نامہ نگاری وہ صحافت ہے جو کسی ضلع، تحصیل یا قصبے کی
سطح پر کی جاتی ہے اور مقامی مسائل کو قومی سطح تک پہنچاتی ہے۔ یہی وہ صحافت ہے
جو عام آدمی کے سب سے قریب ہوتی ہے۔

\medskip
\textbf{اصول:}
\begin{itemize}\setlength{\itemsep}{1pt}
  \item \textbf{مقامی اہمیت} --- وہ خبر بھیجیے جو علاقے کے لوگوں پر اثر ڈالتی
        ہو: پانی، بجلی، سڑک، تعلیم، صحت۔
  \item \textbf{صداقت اور تصدیق} --- چھوٹے علاقے میں افواہ تیزی سے پھیلتی ہے؛
        تصدیق کے بغیر کچھ نہ بھیجیے۔
  \item \textbf{غیر جانبداری} --- مقامی سیاست، برادری اور دھڑے بندی سے بالاتر
        رہنا؛ یہی سب سے بڑا امتحان ہے۔
  \item \textbf{ذاتی دشمنی یا دوستی} کو خبر پر اثر انداز نہ ہونے دینا۔
  \item \textbf{اختصار} --- علاقائی خبر کے لیے جگہ محدود ہوتی ہے۔
  \item \textbf{بروقت ترسیل} --- ڈیڈ لائن کی سختی سے پابندی۔
  \item \textbf{قانونی احتیاط} --- ہتکِ عزت سے بچاؤ؛ الزام کو الزام ہی لکھیے۔
  \item \textbf{تعمیری رویہ} --- صرف شکایت نہیں، مسئلے کا حل اور مثبت پہلو بھی۔
\end{itemize}

\medskip
\textbf{مشکلات:} کم معاوضہ، وسائل کی کمی، مقامی بااثر افراد کا دباؤ، تحفظ کا
فقدان اور تربیت کے مواقع کی کمی۔

\jawab{علاقائی نامہ نگاری ضلعی و مقامی سطح کی صحافت ہے۔ اصول: مقامی اہمیت،
صداقت و تصدیق، غیر جانبداری، ذاتی تعلقات سے بالاتری، اختصار، بروقت ترسیل،
قانونی احتیاط اور تعمیری رویہ۔ بڑی مشکلات کم معاوضہ اور مقامی دباؤ ہیں۔}
\end{hal}

%---------------------------------------------------------------------
\section*{خبر نگاری کی ترقی یافتہ صورتیں}
\markboth{ترقی یافتہ صورتیں}{}

\begin{sawal}{سوال \textenglish{10} \ \textnormal{\small(بہار \textenglish{2013}، سوال \textenglish{5})}}
خبر نگاری کی ترقی یافتہ صورتیں بیان کریں۔
\end{sawal}

\begin{hal}
سادہ خبر صرف \emph{کیا ہوا} بتاتی ہے۔ ترقی یافتہ صورتیں \emph{کیوں}،
\emph{کیسے} اور \emph{آگے کیا} تک جاتی ہیں۔

\medskip
\begin{center}
\begin{tabular}{@{}p{3.6cm} p{10.8cm}@{}}
\toprule
تشریحی / تجزیاتی & واقعے کا پس منظر، اسباب اور ممکنہ نتائج بیان کرنا \\
تحقیقی و تفتیشی & چھپائی گئی حقیقت کو محنت سے بے نقاب کرنا \\
ترقیاتی خبر نگاری & تعلیم، صحت، زراعت اور دیہی ترقی کے مسائل و حل \\
خصوصی / پیشہ ورانہ & معیشت، سائنس، ماحول جیسے شعبوں کی ماہرانہ رپورٹنگ \\
ڈیٹا خبر نگاری & اعداد و شمار کے تجزیے سے کہانی نکالنا؛ گراف اور نقشے \\
شہری صحافت & عام شہری کا موبائل سے خبر یا ویڈیو بھیجنا \\
موبائل صحافت & رپورٹر کا صرف موبائل سے فلم بندی، ادارت اور ترسیل \\
کثیر الوسائل & ایک خبر متن، تصویر، ویڈیو اور انفوگرافک کے ساتھ \\
\bottomrule
\end{tabular}
\end{center}

\jawab{ترقی یافتہ صورتوں میں تشریحی و تجزیاتی، تحقیقی و تفتیشی، ترقیاتی،
خصوصی، ڈیٹا خبر نگاری، شہری صحافت، موبائل صحافت اور کثیر الوسائل رپورٹنگ شامل
ہیں۔ یہ سب سادہ خبر سے آگے بڑھ کر \emph{کیوں} اور \emph{کیسے} کا جواب دیتی
ہیں۔}
\end{hal}

\begin{sawal}{سوال \textenglish{11} \ \textnormal{\small(خزاں \textenglish{2013}، سوال \textenglish{6})}}
تجزیاتی خبر نگاری سے کیا مراد ہے؟ نیز تحقیقی اور تفتیشی خبر نگاری پر نوٹ لکھیں۔
\end{sawal}

\begin{hal}
\textbf{تجزیاتی خبر نگاری۔} وہ رپورٹنگ جس میں محض واقعہ نہیں بتایا جاتا بلکہ
اس کا \emph{پس منظر}، \emph{اسباب} اور \emph{ممکنہ نتائج} بھی بیان کیے جاتے
ہیں --- تاکہ قاری واقعے کا مطلب سمجھ سکے۔ مثلاً روپے کی قدر گرنے کی خبر کے
ساتھ یہ بتانا کہ ایسا کیوں ہوا اور مہنگائی پر کیا اثر پڑے گا۔

\smallskip
\emph{اہم فرق:} تجزیہ \textbf{رائے نہیں} ہوتا۔ اداریہ کہتا ہے ''حکومت کو یہ
کرنا چاہیے''؛ تجزیہ کہتا ہے ''ماہرین کے مطابق اس کے یہ نتائج نکل سکتے ہیں''۔

\medskip
\textbf{تحقیقی و تفتیشی خبر نگاری} (\textenglish{Investigative Reporting})۔
وہ رپورٹنگ جس میں صحافی طویل محنت سے ایسی حقیقت سامنے لاتا ہے جسے کوئی طاقتور
فریق \emph{جان بوجھ کر چھپا} رہا ہو --- بدعنوانی، اختیارات کا ناجائز استعمال،
یا عوامی مفاد کے خلاف کوئی معاملہ۔

\smallskip
\textbf{خصوصیات:}
\begin{itemize}\setlength{\itemsep}{1pt}
  \item مہینوں پر محیط محنت اور ادارے کی مالی و قانونی معاونت۔
  \item دستاویزی ثبوت --- صرف بیان کافی نہیں۔
  \item ایک سے زیادہ آزاد ذرائع سے تصدیق۔
  \item ذرائع کا تحفظ اور صحافی کا اپنا تحفظ۔
  \item شائع کرنے سے پہلے متعلقہ فریق کا مؤقف لینا (قانونی طور پر لازم)۔
\end{itemize}

\smallskip
\textbf{خطرات:} ہتکِ عزت کے مقدمات، دباؤ، دھمکیاں --- اسی لیے یہ صحافت کی سب
سے مشکل اور سب سے قابلِ قدر صورت ہے۔

\jawab{تجزیاتی خبر نگاری واقعے کے ساتھ اس کا پس منظر، اسباب اور نتائج بیان
کرتی ہے، مگر رائے نہیں دیتی۔ تحقیقی و تفتیشی خبر نگاری وہ ہے جس میں صحافی
طویل محنت اور دستاویزی ثبوت سے جان بوجھ کر چھپائی گئی حقیقت بے نقاب کرتا ہے؛
اس میں متعدد ذرائع سے تصدیق، فریقِ مخالف کا مؤقف اور قانونی احتیاط لازم ہے۔}
\end{hal}

%---------------------------------------------------------------------
\section*{مکتوب نگاری، خبر رساں ادارے اور قوانین}
\markboth{مکتوب نگاری، خبر رساں ادارے، قوانین}{}

\begin{sawal}{سوال \textenglish{12} \ \textnormal{\small(بہار \textenglish{2013}، سوال \textenglish{7} اور خزاں \textenglish{2013}، سوال \textenglish{5})}}
مکتوب نگاری کیا ہے؟ اس کی ابتدا اور اہمیت و افادیت بیان کریں۔
\end{sawal}

\begin{hal}
\textbf{تعریف۔} مکتوب نگاری صحافت کی وہ صنف ہے جس میں کسی دوسرے شہر یا ملک میں
مقیم نامہ نگار وہاں کے حالات پر باقاعدگی سے ایک ''مکتوب'' (خط) لکھ کر بھیجتا
ہے، جو اخبار میں مقررہ عنوان کے تحت شائع ہوتا ہے --- مثلاً \emph{مکتوبِ لندن}،
\emph{مکتوبِ دہلی}، \emph{مکتوبِ کراچی}۔ اس میں خبر، تجزیہ اور ذاتی مشاہدہ تینوں
ہوتے ہیں، اور اسلوب نیم ادبی ہوتا ہے۔

\medskip
\textbf{ابتدا۔} جب ذرائع ابلاغ محدود تھے اور تار یا ٹیلی فون عام نہ تھا، دور
دراز کے حالات کی واحد صورت یہی تھی کہ وہاں مقیم کوئی صاحبِ قلم خط لکھ کر
بھیجے۔ اردو صحافت کے ابتدائی دور ہی سے یہ صنف رائج رہی اور بعد میں اس نے
باقاعدہ کالم کی شکل اختیار کر لی۔

\medskip
\textbf{اہمیت و افادیت:}
\begin{itemize}\setlength{\itemsep}{1pt}
  \item دور دراز علاقوں اور بیرونِ ملک کے حالات کی تفصیلی تصویر۔
  \item محض واقعہ نہیں، وہاں کا \emph{ماحول، مزاج اور رائے عامہ} بھی سامنے آتی
        ہے --- جو سادہ خبر نہیں دے سکتی۔
  \item بیرونِ ملک مقیم پاکستانیوں کا وطن سے رابطہ۔
  \item مقامی مسائل کو قومی سطح پر اجاگر کرنے کا ذریعہ۔
  \item اردو نثر اور صحافتی ادب کی ترویج۔
  \item باقاعدگی کے سبب قاری کا اخبار سے مستقل تعلق۔
\end{itemize}

\smallskip
\emph{نوٹ:} اسے ''قارئین کے خطوط'' سے الگ رکھیے --- وہ ادارتی صفحے پر عام
قارئین کی آرا ہوتی ہیں، جبکہ مکتوب ادارے کے نامہ نگار کی باقاعدہ تحریر ہے۔

\jawab{مکتوب نگاری وہ صنف ہے جس میں دوسرے شہر یا ملک کا نامہ نگار وہاں کے
حالات پر باقاعدہ خط لکھ کر بھیجتا ہے (مکتوبِ لندن، مکتوبِ دہلی)۔ اس کی ابتدا
اُس دور میں ہوئی جب رابطے کے تیز ذرائع نہ تھے۔ افادیت: دور دراز حالات کی
تفصیلی تصویر، وہاں کے ماحول و رائے عامہ کی عکاسی، بیرونِ ملک پاکستانیوں کا
رابطہ اور اردو نثر کی ترویج۔}
\end{hal}

\begin{sawal}{سوال \textenglish{13} \ \textnormal{\small(بہار \textenglish{2013}، سوال \textenglish{8} اور خزاں \textenglish{2013}، سوال \textenglish{8})}}
خبر رساں اداروں کا تاریخی پس منظر بیان کریں اور بین الاقوامی خبر رساں اداروں پر
تفصیلی نوٹ لکھیں۔
\end{sawal}

\begin{hal}
\textbf{خبر رساں ادارہ کیا ہے؟} وہ ادارہ جو دنیا بھر سے خبریں جمع کر کے
اخبارات، ریڈیو اور ٹی وی چینلوں کو معاوضے پر فراہم کرتا ہے۔ کوئی ایک اخبار ہر
جگہ اپنا نامہ نگار نہیں رکھ سکتا --- یہی خلا خبر رساں ادارے پر کرتے ہیں۔

\medskip
\textbf{تاریخی پس منظر۔} انیسویں صدی میں تار (\textenglish{telegraph}) کی
ایجاد نے خبر کی تیز ترسیل ممکن بنائی، اور اسی سے یہ ادارے وجود میں آئے:

\begin{center}
\begin{tabular}{@{}p{3.0cm} p{1.8cm} p{2.4cm} p{6.6cm}@{}}
\toprule
\textbf{ادارہ} & \textbf{سنہ} & \textbf{ملک} & \textbf{تفصیل} \\
\midrule
\textenglish{Havas} & \textenglish{1835} & فرانس & دنیا کا پہلا خبر رساں ادارہ؛ بعد میں \textenglish{AFP} بنا \\
\textenglish{AP} & \textenglish{1846} & امریکہ & اخبارات کی مشترکہ تعاونی تنظیم \\
\textenglish{Wolff} & \textenglish{1849} & جرمنی & یورپ کے ابتدائی اداروں میں \\
\textenglish{Reuters} & \textenglish{1851} & برطانیہ & پال جولیس رائٹر؛ آج سب سے معروف \\
\textenglish{UPI} & \textenglish{1907} & امریکہ & \textenglish{AP} کا حریف \\
\textenglish{TASS} & \textenglish{1925} & روس & سرکاری ادارہ \\
\textenglish{Xinhua} & \textenglish{1931} & چین & سرکاری ادارہ \\
\bottomrule
\end{tabular}
\end{center}

\medskip
\textbf{پاکستانی ادارے:} \textenglish{APP} (\textenglish{Associated Press of
Pakistan}، \textenglish{1947}، سرکاری)، \textenglish{PPI} (\textenglish{Pakistan
Press International})، \textenglish{NNI} اور \textenglish{Online}۔

\medskip
\textbf{اہمیت:} بین الاقوامی خبروں کی فراہمی، اخبار کے اخراجات میں کمی، خبر کی
رفتار، اور چھوٹے اخبارات کو بھی عالمی خبروں تک رسائی۔

\textbf{تنقید:} چند مغربی اداروں کی اجارہ داری کے سبب ترقی پذیر ممالک کی تصویر
یکطرفہ پیش ہوتی ہے --- یہی \emph{خبروں کے عالمی بہاؤ میں عدم توازن} کی بحث ہے۔

\jawab{خبر رساں ادارے خبریں جمع کر کے ذرائع ابلاغ کو فراہم کرتے ہیں۔ ابتدا
\textenglish{Havas} (\textenglish{1835}، فرانس) سے ہوئی، پھر \textenglish{AP}
(\textenglish{1846})، \textenglish{Reuters} (\textenglish{1851})،
\textenglish{UPI}، \textenglish{TASS} اور \textenglish{Xinhua}۔ پاکستان میں
\textenglish{APP} اور \textenglish{PPI} اہم ہیں۔ ان سے رفتار اور رسائی بڑھی،
مگر مغربی اجارہ داری پر تنقید بھی ہوتی ہے۔}
\end{hal}

\begin{sawal}{سوال \textenglish{14} \ \textnormal{\small(خزاں \textenglish{2013}، سوال \textenglish{2})}}
قوانینِ صحافت پر تفصیلی نوٹ لکھیں۔
\end{sawal}

\begin{hal}
\textbf{آئینی بنیاد۔} آئینِ پاکستان \textenglish{1973} کا \textbf{آرٹیکل
\textenglish{19}} اظہارِ رائے اور آزادیٔ صحافت کی ضمانت دیتا ہے، مگر ''معقول
پابندیوں'' کے ساتھ: اسلام کی عظمت، ملکی سلامتی، امنِ عامہ، اخلاق، اور عدالت کی
توہین۔ \textbf{آرٹیکل \textenglish{19-A}} (\textenglish{2010}) معلومات تک
رسائی کا حق دیتا ہے۔

\medskip
\begin{center}
\begin{tabular}{@{}p{6.4cm} p{8.0cm}@{}}
\toprule
\textenglish{Press and Publications Ordinance 1963} & سخت ترین دور؛ \textenglish{1988} میں منسوخ \\
\textenglish{RPPPO 2002} & پریس اور اشاعت کی رجسٹریشن \\
\textenglish{PEMRA Ordinance 2002} & نجی ریڈیو و ٹی وی کا ضابطہ \\
\textenglish{Press Council of Pakistan Ordinance 2002} & ضابطۂ اخلاق اور شکایات \\
\textenglish{Defamation Ordinance 2002} & ہتکِ عزت \\
\textenglish{Freedom of Information Ordinance 2002} & سرکاری معلومات تک رسائی \\
\textenglish{PECA 2016} & برقی جرائم اور آن لائن مواد \\
\textenglish{Right of Access to Information Act 2017} & معلومات کا حق \\
\bottomrule
\end{tabular}
\end{center}

\medskip
\textbf{رپورٹر کے لیے عملی احتیاطیں:}
\begin{itemize}\setlength{\itemsep}{1pt}
  \item \textbf{ہتکِ عزت} --- الزام کو ہمیشہ الزام لکھیے، ثابت شدہ حقیقت نہیں۔
  \item \textbf{توہینِ عدالت} --- زیرِ سماعت مقدمے پر رائے زنی نہ کیجیے۔
  \item \textbf{ملزم اور مجرم کا فرق} --- عدالت کے فیصلے سے پہلے ہر شخص ملزم
        ہے۔
  \item متاثرین، خاص طور پر بچوں اور خواتین کی شناخت ظاہر نہ کیجیے۔
\end{itemize}

\jawab{بنیاد آرٹیکل \textenglish{19} اور \textenglish{19-A} ہیں۔ اہم قوانین:
\textenglish{PPO 1963} (منسوخ)، \textenglish{RPPPO 2002}،
\textenglish{PEMRA 2002}، پریس کونسل آرڈیننس، ہتکِ عزت آرڈیننس
\textenglish{2002}، \textenglish{PECA 2016} اور معلومات تک رسائی کا ایکٹ
\textenglish{2017}۔ رپورٹر کے لیے ہتکِ عزت، توہینِ عدالت اور ملزم/مجرم کا فرق
سب سے اہم احتیاطیں ہیں۔}
\end{hal}
